\documentclass[11pt]{article}

% -------------------------------------------------
% Packages
% -------------------------------------------------
\usepackage[utf8]{inputenc}
\usepackage[T1]{fontenc}
\usepackage{amsmath, amssymb, mathtools}
\usepackage{geometry}
\usepackage{booktabs}
\usepackage{hyperref}
\usepackage{enumitem}

\geometry{margin=1in}

% -------------------------------------------------
% Document
% -------------------------------------------------
\begin{document}

\begin{center}
    {\LARGE \textbf{Exercise 2.2 — Transition Matrix from Conditional Expectations}}\\[0.5em]
    {\large Recursive Macroeconomic Theory}\\[0.25em]
    {\normalsize Lars Ljungqvist and Thomas J. Sargent}
\end{center}

\vspace{1em}

\section*{Problem}

Consider a two-state Markov chain and a random variable \(y_t=\bar{y}^{\top}x_t\) where

\[
\bar y =
\begin{pmatrix}
1 \\
5
\end{pmatrix}.
\]
It is known that \(E(y_{t+1}|x_t) = \begin{pmatrix}
    1.8 \\
    3.4
\end{pmatrix}\) and that \(E(y_{t+1}^2|x_t) = \begin{pmatrix}
    5.8\\
    15.4
\end{pmatrix}.\)

We find a transition matrix consistent with these conditional expectations and determine whether the matrix is unique.


\section*{Analytical Solution}

For a two-state Markov chain, the transition matrix is 

\[
P = \begin{pmatrix}
    p_{00} & p_{01} \\
    p_{10} & p_{11}
\end{pmatrix}
\]
By definition of \(P_{ij}= Pr(x_{t+1}=e_j | x_t = e_i)\), we know that \(p_{00} +p_{01}=1\) and that \(p_{10} + p_{11}= 1\).

Furthermore, we know that the conditional expectation of a function of next period's state is 
\[
E(y_{t+1}|x_t)= P\bar{y}.
\]
This implies we have to solve the following system in matrix notation:

\[
\begin{pmatrix}
    p_{00} & p_{01} \\
    p_{10} & p_{11}
\end{pmatrix}
\begin{pmatrix}
    1 \\
    5
\end{pmatrix}=\begin{pmatrix}
    1.8 \\
    3.4
\end{pmatrix},
\]
which leads to the following transition matrix:

\[
P=
\begin{pmatrix}
    0.8 & 0.2 \\
    0.4 & 0.6
\end{pmatrix}.
\]
Note that, since the two state values are distinct, the row-sum restriction together with
the conditional mean provides two linearly independent equations for the two
unknown probabilities in each row. Therefore, the transition matrix is uniquely
determined.

Moreover, since \(P_{ij}>0\) for all \(i,j\), Theorem 2.2.1 implies that the
Markov chain has a unique stationary distribution and is asymptotically stationary.

Thus, we are only left with checking the second moment condition

\[
y^2 = \begin{pmatrix}
    1^2 \\
    5^2
\end{pmatrix}= \begin{pmatrix}
    1 \\
    25
\end{pmatrix}.
\]
Therefore,

\[
P \begin{pmatrix}
    1 \\
    25
\end{pmatrix}= \begin{pmatrix}
    0.8 + 5\\
    0.4 + 15
\end{pmatrix} = \begin{pmatrix}
    5.8 \\
    15.4
\end{pmatrix}= E(y_{t+1}^2|x_t)
\]


\section*{Two-State Business-Cycle Model Extension}

In Exercise 2.1, the transition matrix was treated as known and used to study
the properties of the Markov chain. A natural econometric extension reverses
the problem: instead of taking \(P\) as given, we seek to recover it from
observable implications of the stochastic process. In empirical applications, however, conditional moments are estimated from
finite samples and therefore need not satisfy the theoretical restrictions of
the model exactly. This motivates a minimum-distance approach.

In particular, think of the economy as being in either a recessionary regime or an expansionary regime, with \(y_t\) representing quarterly real GDP growth.

Suppose we approximate GDP growth by two representative values:
\[
\bar{y} =\begin{pmatrix}
    -2 \\
    3
\end{pmatrix}, 
\]
so state 0 means recession and state 1 means expansion. To be clear, these numbers are just illustrative rather than estimates from a particular dataset.

Suppose that, in an empirical application, we obtained the following sample
estimates of the conditional moments:
\[
E(y_{t+1}|x_t)=\begin{pmatrix}
    -0.8 \\
    2.4
\end{pmatrix} \quad \text{and} \quad E(y_{t+1}^2|x_t)=\begin{pmatrix}
    5.4 \\
    8.1
    \end{pmatrix}.
\]
Of course, these sample estimates may not satisfy the theoretical restrictions imposed by the Markov model exactly; this is prevented by sample variation. Thus, instead of demanding a P satisfying every equation exactly, we can look for the transition matrix that comes closest to matching all of the moments simultaneously.

For each current state \(i\), we write
\[
q_i=P(x_{t+1}= \text{expansion}| x_t =i).
\]
Then, 

\[P_{i,.} = (1-q_i, q_i), \]

and, for any proposed \(q_i\), the model predicts 
\[
E(y_{t+1}|x_t)= (1-q_i)(-2) +q_i(3) \quad \text{and} \quad E(y_{t+1}^2|x_t)=(1-q_i)4+q_i9.
\]
We choose \(q_i\) to minimize 

\[
(E_{model}(y_{t+1}|x_t = i)-E_{data}(y_{t+1}|x_t=i))^2+(E_{model}(y_{t+1}^2|x_t = i)-E_{data}(y_{t+1}^2|x_t=i))^2,
\]
which is a simple minimum-distance estimator. In the Python companion of this exercise, we estimate the transition matrix to be

\[
\hat{P}=\begin{pmatrix}
    0.74 & 0.26 \\
    0.15 & 0.85
\end{pmatrix}
\]

The estimated matrix implies that, conditional on being in a recession, the
economy remains in recession with probability \(0.74\) and moves into expansion
with probability \(0.26\). Conditional on being in an expansion, the expansion
persists with probability \(0.85\). Unlike the exact solution to Exercise 2.2,
\(\hat P\) does not reproduce every sample moment exactly; rather, it provides
the transition matrix that best reconciles the two sets of noisy conditional
moments according to the chosen minimum-distance criterion.

\section*{Reference}

Ljungqvist, L. and Sargent, T. J. (2018).
\textit{Recursive Macroeconomic Theory},
4th edition, MIT Press.

\end{document}
